% ===================================================================
%  جدول نمبر 6 — گھر کے اندر، سامان، خوراک، روشنی۔
%
%  Traduction, faite en 2026, en ourdou d'aujourd'hui, sur l'ido de
%  texto/io. Regles de la colonne : voir 10-jadval-01.tex. Cinq
%  pieces, cinq numerotations : les cles portent c1 a c5. Ecrit avec
%  outils/ordo.py sous les yeux.
%
%  CINQUANTE-DEUX PAIRES, ET PRESQUE TOUTES SE RETOURNENT. Ce tableau
%  ne fait rien d'autre que placer des objets les uns par rapport aux
%  autres : le seau DEVANT la table de toilette, l'armoire QUI
%  contient le linge, l'edredon SUR le couvre-lit, le cadre PENDU au
%  mur, la marmite AVEC son couvercle, le tonneau MUNI de son
%  robinet. Chaque preposition de l'ido est une postposition en
%  ourdou, et une postposition suit ce qu'elle regit : l'ordre
%  s'inverse a chaque fois. Le tableau 5, deux fois plus long, n'en
%  retournait que onze.
%
%  LA MAISON DU LIVRET N'EST PAS UNE MAISON DU PENDJAB, ET LA COLONNE
%  NE LA DEPLACE PAS. Une maison de 1926 a Lahore n'a ni lit a
%  baldaquin, ni cheminee a chenets, ni baignoire : on y dort sur la
%  چارپائی, on s'y chauffe a l'انگیٹھی qu'on porte de piece en piece,
%  on s'y lave au لوٹا. La tentation etait d'ecrire چارپائی pour le
%  lit — mot vivant, connu de tous, et faux. La colonne ecrit پلنگ,
%  qui est le lit a montants, et آتش دان la cheminee de chambre,
%  et غسل کا ٹب la baignoire.
%
%  MAIS انگیٹھی EST GARDE AU BLOC c4-01-1, comme le 아궁이 coreen et
%  le அடுப்பு tamoul au meme bloc. La, la chose EST un foyer ouvert
%  ou tourne une broche, et انگیٹھی la designe exactement. Ailleurs
%  c'est un autre objet et un autre mot. Trois colonnes sans parente
%  ont fait le meme partage au meme endroit.
%
%  LA CUISINE EST L'ENDROIT OU L'OURDOU EST LE PLUS ENTIER, ET IL
%  L'EST EN DEUX LANGUES A LA FOIS : چمٹا les pincettes, بھٹی le
%  foyer, دھونکنی le soufflet, کوئلہ le charbon, چھلنی le tamis,
%  کفگیر l'ecumoire, کڑاہی la poele, دیگچی la casserole, ہانڈی le
%  pot, ڈھکن le couvercle, جھاڑو le balai, قیف l'entonnoir, ٹونٹی le
%  robinet, راکھ la cendre, لکڑی la buche. Quinze objets, treize mots
%  indiens et deux arabes — et l'on ne saurait dire, a l'usage,
%  lequel est lequel.
%
%  LE FAC-SIMILE REEMPLOIE (6) — le savon au bloc c1-02-1, la femme
%  de chambre au bloc c1-06-1 — et la source ne bouge pas. Il donne
%  aussi DEUX notes sur le feuillet 41, chacune avec son propre
%  asterisque remis a un : deux blocs, pas un.
% ===================================================================

%%K t06-apar-1 apar
\VUsaut{6.0mm}
\VUtitre{15.0pt}{60}{جدول نمبر 6}
\VUsaut{1.4mm}
\VUtitre{11.4pt}{0}{گھر کے اندر، سامان، خوراک،}
\VUsaut{0.4mm}
\VUtitre{11.4pt}{0}{روشنی۔}
\VUsaut{2.2mm}

%%K t06-apar-2 apar
گھر مکمل ہو گیا۔ یوآنیس کے ماموں اپنے نئے مکان میں آباد ہو گئے ہیں ؛
اُس کی ترتیب ہم جانتے ہیں۔ اب دیکھیں کہ اُنھوں نے اپنا گھر کیسے سجایا
ہے۔ پہلے یہ دیکھتے ہیں :

%%K t06-tit-1 sub
\VUpk{10.2pt}{40}{سونے کا کمرہ۔}
\VUsaut{3.0mm}

%%K t06-c1-01-1 p
1. --- ہم بے وقت آ گئے، کیونکہ خادمہ \VUgras{کمرہ} صاف کرنے میں لگی
ہے۔

%%K t06-c1-02-1 p
2. --- \VUgras{بناؤ سنگھار کی میز} \textsuperscript{(1)} پر ادھر ادھر
کئی چیزیں ہیں ؛ اُن سے منہ ہاتھ دھوتے، کنگھی کرتے، مختصراً بناؤ
سنگھار کرتے ہیں۔ وہ \VUgras{چھوٹا طشت} \textsuperscript{(2)} اور
\VUgras{پانی کا جگ} \textsuperscript{(3)}، \VUgras{بالوں کا برش}
\textsuperscript{(4)}، \VUgras{کنگھا} \textsuperscript{(5)}،
\VUgras{صابن} \textsuperscript{(6)}، \VUgras{اسفنج}
\textsuperscript{(7)}، اور آخر میں دانت صاف کرنے کے مائع کے لیے
\VUgras{چھوٹی شیشی} \textsuperscript{(8)} ہیں۔

%%K t06-c1-03-1 p
3. --- فرش پر، بناؤ سنگھار کی میز کے سامنے، یہ رہی گندا پانی جمع کرنے
کی \VUgras{بالٹی} \textsuperscript{(9)}۔

%%K t06-c1-04-1 p
4. --- \VUgras{ڈیوڑھی} \textsuperscript{(10)} میں ہمیں کچھ
\VUgras{کپڑوں کی کھونٹیاں} \textsuperscript{(13)} اور دو
\VUgras{چھتریاں} \textsuperscript{(11)} نظر آتی ہیں ؛ وہ چھتریاں
\VUgras{چھتری دان} \textsuperscript{(12)} میں ہیں۔

%%K t06-c1-05-1 p
5. --- دروازے کی دوسری طرف \VUgras{شیشے والی الماری}
\textsuperscript{(14)} ہے ؛ اُس میں \VUgras{کپڑے} \textsuperscript{(15)}
ہیں۔

%%K t06-c1-06-1 p
6. --- \VUgras{خادمہ} \textsuperscript{(6)} جس وقت \VUgras{پلنگ}
\textsuperscript{(17)} درست کر رہی ہے، اُس وقت اُس کے مختلف حصے دیکھ
لیں۔ \VUgras{پلنگ کے ڈھانچے} \textsuperscript{(20)} میں اچھا
\VUgras{کمانی دار گدا} \textsuperscript{(21)} ہے ؛ اُس پر یوستینا ابھی
\VUgras{اونی گدا} \textsuperscript{(22)} بچھائے گی۔ پھر کرسیوں سے دو
\VUgras{چادریں} \textsuperscript{(23)} اور \VUgras{کمبل}
\textsuperscript{(24)} اٹھائے گی۔ تب پلنگ کے سرہانے
\VUgras{لمبا تکیہ} \textsuperscript{(25)} اور \VUgras{تکیہ}
\textsuperscript{(26)} رکھے گی ؛ وہ \VUgras{پروں والی رضائی}
\textsuperscript{(27)} کے ساتھ \VUgras{پلنگ پوش} \textsuperscript{(32)}
پر ہیں۔ \VUgras{پردوں} \textsuperscript{(19)} اور \VUgras{پلنگ کی
چاندنی} \textsuperscript{(18)} کی گرد جھاڑنے کے لیے وہ پروں کا جھاڑن
استعمال کرتی ہے ؛ یوں کام کا ایک حصہ مکمل ہو جاتا ہے۔

%%K t06-c1-07-1 p
7. --- پھر اُسے \VUgras{پلنگ کے پاس کی میز} \textsuperscript{(28)}
کچھ ٹھیک کرنی ہو گی، \VUgras{جگانے والی گھڑی} \textsuperscript{(29)}
کو چابی دینی ہو گی، \VUgras{رات کا چراغ} \textsuperscript{(30)} تیار
کرنا ہو گا، اور شمع دان صاف کرنے کے لیے \VUgras{موم بتی}
\textsuperscript{(31)} اٹھانی ہو گی۔

%%K t06-tit-2 sub
\VUsaut{5.0mm}
\VUpk{10.2pt}{40}{سلائی کی ٹوکری اور انگیٹھی کا سامان۔}
\VUsaut{3.0mm}

%%K t06-c1-08-1 p
8. --- \VUgras{سلائی کی ٹوکری} \textsuperscript{(34)} — وہ
\VUgras{چھوٹی چوکی} \textsuperscript{(33)} کے پاس ہے — بھی بہت ترتیب
مانگتی ہے۔ اُسے \VUgras{کڑھائی کے کام} \textsuperscript{(35)} تہ کرنے
ہوں گے ؛ \VUgras{سلائی کی میز} \textsuperscript{(38)} میں
\VUgras{انگشتانہ}، \VUgras{دھاگا} \textsuperscript{(36)}،
\VUgras{سوئیاں} \textsuperscript{(37)} اور \VUgras{قینچی}
\textsuperscript{(39)} رکھنے ہوں گے ؛ اور میز کا ڈھکن
\VUgras{چابی} \textsuperscript{(40)} سے بند کرنا ہو گا۔

%%K t06-c1-09-1 p
9. --- بیچ میں رکھی \VUgras{ایک پائے والی میز} \textsuperscript{(41)}
پر یوستینا \VUgras{صراحی} \textsuperscript{(42)} کا پانی بدلے گی۔
\VUgras{چینی} \VUgras{چینی دان} \textsuperscript{(43)} میں ڈالے گی،
\VUgras{چینی کی چمٹی} \textsuperscript{(44)} صاف کرے گی، اور
\VUgras{کپڑوں کا برش} \textsuperscript{(46)} اٹھا لے گی۔

%%K t06-c1-10-1 p
10. --- کمرے کی دائیں طرف اُس سے کم کام مانگے گی ؛ کیونکہ
\VUgras{لمبی آرام کرسی} \textsuperscript{(47)} اور اُس کی
\VUgras{گدی} \textsuperscript{(48)} اپنی ٹھیک جگہ پر ہیں ؛ اور سردی نہ
ہونے کی وجہ سے \VUgras{آتش دان} \textsuperscript{(49)} کے سامان میں سے
کوئی چیز اپنی جگہ سے نہیں ہٹی۔ \VUgras{راکھ کا بیلچہ}
\textsuperscript{(50)}، \VUgras{بڑا چمٹا} \textsuperscript{(51)}،
\VUgras{لکڑی کے سہارے} \textsuperscript{(55)}، \VUgras{راکھ کی روک}
\textsuperscript{(56)} صاف ہیں ؛ \VUgras{آگ کا پردہ}
\textsuperscript{(54)} اپنی جگہ سے نہیں ہٹا ؛ اور \VUgras{لکڑی کا
صندوق} \textsuperscript{(52)} \VUgras{لکڑیوں} \textsuperscript{(53)} سے
اچھی طرح بھرا ہے۔

%%K t06-noto-1 noto
(*) Babo = ننھا بچہ ؛ ا۔ baby، ف۔ bébé، ا۔ bambino۔

%%K t06-noto-2 noto
(*) Lambrequino = ف۔ lambrequin، ہ۔ lambrequines، پ۔ lambrequins ؛
ج۔ ا۔ میں بھی lambrequins ؛ پس ج۔ ا۔ ف۔ پ۔ ہ۔

%%K t06-c1-11-1 p
11. --- پھر بھی اُسے \VUgras{پنڈولم گھڑی} \textsuperscript{(57)}،
\VUgras{شمع دانوں} \textsuperscript{(58)} اور \VUgras{چھوٹی
ڈبیوں} \textsuperscript{(59)} کی گرد جھاڑنی ہو گی ؛ آخر میں
\VUgras{آئینہ} \textsuperscript{(60)} پونچھنا ہو گا۔

%%K t06-c1-12-1 p
12. --- ننھے بچے (بابو (*)) کا \VUgras{جھولا} \textsuperscript{(61)}
پہلے ہی تیار ہے۔

%%K t06-tit-3 sub
\VUsaut{5.0mm}
\VUpk{10.2pt}{40}{بیٹھک۔}
\VUsaut{3.0mm}

%%K t06-c2-01-1 p
1. --- اب یہ رہی بیٹھک۔ وہ بہت خوبصورت \VUgras{کمرہ} ہے۔ دائیں کونے
میں رکھا آتش دان نفیس \VUgras{چھوٹے مجسمے} \textsuperscript{(1)} اور
دو خوبصورت \VUgras{گلدانوں} \textsuperscript{(3)} سے سجا ہے ؛ اور
مخملی \VUgras{جھالر} \textsuperscript{(2)} (*) سے ڈھکا ہے۔

%%K t06-c2-02-1 p
2. --- چولھے کی چنگاریاں قالین کو جلا نہ دیں، اِس لیے چولھے کے سامنے
تانبے کی \VUgras{چنگاری روک} \textsuperscript{(4)} رکھی ہے۔

%%K t06-c2-03-1 p
3. --- قریب ہی خواتین کی نفیس \VUgras{لکھنے کی میز}
\textsuperscript{(5)} کھلی ہے۔ اُس پر \VUgras{گھر کی مالکن}
\textsuperscript{(23)} اپنے خط لکھتی ہے۔

%%K t06-c2-04-1 p
4. --- کھڑکی پر \VUgras{باریک پردے} \textsuperscript{(6)} ہیں ؛ اُن کے
\VUgras{بند} \textsuperscript{(8)} ہیں، جو \VUgras{کھونٹیوں}
\textsuperscript{(7)} میں اٹکے ہیں ؛ اور \VUgras{اوپر سے جھالر دار}
\textsuperscript{(9)} موٹے پردے بھی ہیں۔

%%K t06-c2-05-1 p
5. --- کھڑکی کے طاق میں ایک \VUgras{گملے کا خول}
\textsuperscript{(11)} ہے ؛ اُس میں سبز \VUgras{پودا}
\textsuperscript{(10)} ہے — شاندار کھجور نما پودا ؛ اُس کے پتے
تقریباً \VUgras{مٹی کے تیل کے لیمپ} \textsuperscript{(12)} تک پھیلے
ہیں۔

%%K t06-c2-06-1 p
6. --- لیمپ کا \VUgras{روشنی کا غلاف} \textsuperscript{(13)} ماریا نے
سجایا ہے — وہ دلکش \VUgras{نوجوان لڑکی} \textsuperscript{(14)} ؛ وہ
پیانو \textsuperscript{(15)} کے پاس ہے۔ \VUgras{چوکی}
\textsuperscript{(16)} پر بہت سیدھی بیٹھ کر ایک موسیقی کا سبق سیکھ رہی
ہے۔ وہ بہت ماہر اور محتاط ہے ؛ اُس کی \VUgras{موسیقی کی الماری}
\textsuperscript{(17)} اِس کا ثبوت ہے — وہ بالکل ترتیب سے ہے۔

%%K t06-tit-4 sub
\VUsaut{5.0mm}
\VUpk{10.2pt}{40}{شرارتی لڑکا۔}
\VUsaut{3.0mm}

%%K t06-c2-07-1 p
7. --- اُس کا بھائی پالس ایک \VUgras{کتاب} \textsuperscript{(19)} ڈھونڈ
رہا ہے ؛ وہ \VUgras{کتابوں کی الماری} \textsuperscript{(18)} میں ہے۔
وہ \VUgras{نوجوان} \textsuperscript{(20)} ہے، بے چین اور ناقابلِ
برداشت۔ بہت اوپر رکھی کتاب تک پہنچنے کے لیے الماری سے لٹک رہا ہے ؛
اوپر رکھے \VUgras{نیم تن مجسمے} \textsuperscript{(21)} کو گرا دینے کا
خطرہ ہے۔

%%K t06-c2-08-1 p
8. --- یہ حماقت کرنے کے لیے وہ ایک موقع سے فائدہ اٹھا رہا ہے : اُس کی
ماں ایک سہیلی سے باتوں میں لگی ہے۔ وہ \VUgras{خاتون}
\textsuperscript{(22)} کچھ دن اُن کے گھر رہنے آئی ہے۔ ماں
\VUgras{صوفے} \textsuperscript{(24)} پر بیٹھی ہے ؛ وہ \VUgras{آرام
کرسی} \textsuperscript{(25)} کے پاس ہے۔ اُس کی سہیلی \VUgras{کرسی}
\textsuperscript{(27)} پر ہے ؛ ہمیں اُس کی صرف \VUgras{پشت}
\textsuperscript{(26)} نظر آتی ہے۔

%%K t06-c2-09-1 p
9. --- دیوار پر لٹکے بڑے \VUgras{چوکھٹے} \textsuperscript{(30)} میں
ایک رشتے دار کی \VUgras{تصویر} \textsuperscript{(29)} ہے۔ میز پر رکھے
\VUgras{تصویروں کے مرقع} \textsuperscript{(32)} میں باقی گھر والوں کی
\VUgras{تصویریں} \textsuperscript{(33)} جمع ہیں۔ بچوں نے ابھی اُسے
ورق گردانی کی ہے ؛ کیونکہ \VUgras{کرسیاں} \textsuperscript{(31)} اب
تک میز کے گرد جمع ہیں۔

%%K t06-c2-10-1 p
10. --- چینی مٹی کا \VUgras{چینی گلدان} \textsuperscript{(34)} — وہ
\VUgras{کاغذ کاٹنے کی چھری} \textsuperscript{(35)} کے پاس ہے — ماریا
کو اُس کے تہوار پر تحفے میں ملا تھا ؛ اور کمرے کا کونہ سجانے والا
\VUgras{اوٹ کا پردہ} \textsuperscript{(36)} اُس کے ماموں چین سے لائے
تھے۔

%%K t06-c2-11-1 p
11. --- بہت خوبصورت \VUgras{تصویریں} \textsuperscript{(37)} اُسے
خوشگوار بناتی ہیں ؛ وہ \VUgras{دیوار} \textsuperscript{(42)} پر لٹکی
ہیں۔ \VUgras{چھت کا فانوس} \textsuperscript{(38)} اُسے روشنی دیتا
ہے ؛ اُس کے \VUgras{بجلی کے قمقمے} \textsuperscript{(39)} شام کو
خوشی سے چمکتے ہیں۔ اِس لیے یہ بیٹھک بہت خوشگوار لگتی ہے۔ لکڑی کے
فرش پر بچھا نرم \VUgras{قالین} \textsuperscript{(40)} اور بہت آسانی
سے کام کرنے والی \VUgras{بجلی کی گھنٹی} \textsuperscript{(41)} اُس کے
آرام کو پورا کرتے ہیں۔

%%K t06-tit-5 sub
\VUsaut{3.6mm}
\VUpk{10.2pt}{40}{کھانے کا کمرہ۔}
\VUsaut{3.0mm}

%%K t06-c3-01-1 p
1. --- رات کے کھانے کی گھنٹی بجی۔ ماریا کے سوا سارا خاندان میز کے گرد
جمع ہے۔ چینی مٹی کی بہترین \VUgras{انگیٹھی} \textsuperscript{(1)}
کھانے کے کمرے کو خوشگوار طور پر گرم رکھتی ہے ؛ اُس کی باریک
\VUgras{نلی} \textsuperscript{(2)} ہے۔

%%K t06-c3-02-1 p
2. --- اِس کمرے میں زیادہ سامان نہیں۔ دیوار کے ساتھ، \VUgras{چھوٹی
چوکی} \textsuperscript{(4)} کے اوپر ایک سجاوٹی \VUgras{چھوٹا فوارہ}
\textsuperscript{(3)} لٹکا ہے۔ بالکل قریب \VUgras{پیش خوان}
\textsuperscript{(5)} ہے ؛ اُس پر ٹھنڈے کھانے، میٹھے کی رکابیاں، اور
میز کی وہ سب چیزیں رکھی جاتی ہیں جن کی فوراً ضرورت نہیں۔

%%K t06-c3-03-1 p
3. --- یوں اوپر والے تختے پر ہمیں \VUgras{بھنا ہوا مرغی کا بچہ}
\textsuperscript{(7)}، \VUgras{مچھلی} \textsuperscript{(8)} اور
\VUgras{پیالہ} \textsuperscript{(10)} نظر آتے ہیں — اُس میں
\VUgras{پھل} \textsuperscript{(9)} ہیں۔ نیچے یہ رہے \VUgras{پنیر}
\textsuperscript{(11)}، \VUgras{روٹی} \textsuperscript{(12)}،
\VUgras{شیشیوں کی ٹرے} \textsuperscript{(13)} — اُس میں سلاد کو مزے
دار بنانے کی سب چیزیں ہیں : تیل، سرکہ، \VUgras{نمک}
\textsuperscript{(14)}، \VUgras{کالی مرچ} \textsuperscript{(15)}۔ آخر
میں نیچے والے تختے پر \VUgras{کیک} \textsuperscript{(17)} ہے ؛ وہ
\VUgras{رائی دان} \textsuperscript{(16)} کے پاس ہے۔

%%K t06-c3-04-1 p
4. --- \VUgras{دیوار گیر گھڑی} \textsuperscript{(20)} کہلانے والی
کھانے کے کمرے کی گھڑی کی شکل بہت انوکھی ہے۔ وہ لکڑی کے چوکھٹے میں
جڑی ہے ؛ اُس چوکھٹے میں \VUgras{ہوا کا دباؤ ناپنے والا}
\textsuperscript{(19)} اور \VUgras{حرارت پیما} \textsuperscript{(18)}
بھی ہیں۔

%%K t06-tit-6 sub
\VUsaut{4.6mm}
\VUpk{10.2pt}{40}{رات کا کھانا پیش کرنا۔}
\VUsaut{3.0mm}

%%K t06-c3-05-1 p
5. --- کمرے کی سب دیواروں پر موم کیے بلوط کی \VUgras{لکڑی کی
تختیاں} \textsuperscript{(21)} لگی ہیں۔ برآمدے کی طرف کھلنے والے
دروازے کو \VUgras{دروازے کا پردہ} \textsuperscript{(22)} اچھی طرح بند
کرتا ہے۔ رات کے کھانے کے بعد خاندان وہاں کافی پینے جائے گا۔
\VUgras{پیالیاں} \textsuperscript{(24)} اور \VUgras{طشتریاں}
\textsuperscript{(25)} پہلے ہی \VUgras{رکابیوں کی الماری}
\textsuperscript{(23)} پر تیار ہیں۔

%%K t06-c3-06-1 p
6. --- میز کے اوپر چھت میں \VUgras{لٹکتا چراغ} \textsuperscript{(26)}
لگا ہے ؛ اُس کا بہت چمکدار قمقمہ شام کو سارے کھانے کے کمرے کو روشن
کرتا ہے۔

%%K t06-c3-07-1 p
7. --- اب خود \VUgras{کھانے کی میز} \textsuperscript{(27)} دیکھیں۔ جب
خاندان اندر آیا تو میز تیار تھی۔ \VUgras{میز پوش}
\textsuperscript{(28)} پر، ہر کھانے والے کی جگہ \VUgras{رکابی}
\textsuperscript{(33)}، \VUgras{منہ کا رومال} \textsuperscript{(29)}،
\VUgras{چمچہ} \textsuperscript{(34)}، \VUgras{کانٹا}
\textsuperscript{(35)}، \VUgras{چھری} \textsuperscript{(36)} اور
\VUgras{گلاس} \textsuperscript{(32)} رکھے تھے۔ شراب کے لیے
\VUgras{بوتلیں} \textsuperscript{(30)} اور پانی کے لیے
\VUgras{صراحی} \textsuperscript{(31)} بھی تھیں۔

%%K t06-c3-08-1 p
8. --- \VUgras{خادم} \textsuperscript{(39)} پہلے \VUgras{شوربے کا
پیالہ} \textsuperscript{(37)} لایا ؛ اُس میں کچھ شوربہ اب تک بھاپ
اڑا رہا ہے۔ اب وہ پہلا \VUgras{کھانا} \textsuperscript{(38)} — گوشت
کا کھانا — پیش کر رہا ہے۔ پھر وہ چیزیں لائے گا جن کے نام ہم لے چکے
ہیں ؛ آخر میں \VUgras{میٹھے} کے بعد، جب اُس کے مالک برآمدے میں چلے
جائیں گے، وہ میز کا سامان اٹھا کر کھانے کا کمرہ پھر ترتیب دے گا۔

%%K t06-c3-08-2 p
\textit{نوٹ۔} --- \textit{خوراک سے متعلق عام الفاظ یہاں بہت}
\textit{مختصراً دیے گئے ہیں۔ « ہوٹل » (نمبر 13) اور « بازار »}
\textit{(نمبر 15) کے دو جدولوں میں یہ فہرست مکمل کی جائے گی۔}

%%K t06-tit-7 sub
\VUsaut{4.6mm}
\VUpk{10.2pt}{40}{باورچی خانہ۔}
\VUsaut{3.0mm}

%%K t06-c4-01-1 p
1. --- آدھے کھلے دروازے سے جس باورچی خانے میں ہم جھانکتے ہیں، وہاں
خوبصورت بے ترتیبی نظر آتی ہے۔ \VUgras{انگیٹھی} \textsuperscript{(1)}
کے سامنے — اُس میں \VUgras{آگ} \textsuperscript{(3)} چٹخ رہی ہے —
\VUgras{سیخ گھمانے کا آلہ} \textsuperscript{(5)} لگا ہے۔ خوبصورت
\VUgras{بھنا گوشت} \textsuperscript{(7)} \VUgras{سیخ میں پرویا}
\textsuperscript{(6)} ہوا پک رہا ہے۔

%%K t06-c4-02-1 p
2. --- ابھی استعمال کی گئی \VUgras{دھونکنی} \textsuperscript{(8)} اور
\VUgras{کوئلے کا بیلچہ} \textsuperscript{(9)} \VUgras{لکڑیوں}
\textsuperscript{(10)} کی طرح فرش پر ہی رہ گئے۔

%%K t06-c4-03-1 p
3. --- انگیٹھی کا اوپر کا حصہ ترتیب سے ہے۔ \VUgras{لالٹین}
\textsuperscript{(11)}، \VUgras{ماچس کی ڈبیا} \textsuperscript{(13)} —
اُس میں \VUgras{ماچس کی تیلیاں} \textsuperscript{(12)} ہیں —،
\VUgras{شمع دان} \textsuperscript{(14)}، اور \VUgras{موم بتی کا
غلاف} \textsuperscript{(15)} — اُس میں \VUgras{موم بتی}
\textsuperscript{(16)} ہے — اپنی جگہ پر ہیں۔ مگر بڑی \VUgras{کوئلے
کی بالٹی} \textsuperscript{(18)} — وہ \VUgras{پتھر کے کوئلے}
\textsuperscript{(17)} سے بھری ہے، جسے \VUgras{باورچن}
\textsuperscript{(23)} نے ابھی تہ خانے میں بھرا — باورچی خانے کے
بیچ ہی میں رہ گئی۔

%%K t06-c4-04-1 p
4. --- سچ تو یہ ہے کہ بے چاری عورت کو جلدی کرنی ہے تاکہ ناشتہ ٹھیک
وقت پر تیار ہو۔ اب وہ \VUgras{چولھے} \textsuperscript{(19)} کے پاس ہے
اور \VUgras{شوربہ} \textsuperscript{(24)} ہلا رہی ہے ؛ وہ زیادہ جل نہ
جائے۔

%%K t06-c4-05-1 p
5. --- \VUgras{تنور} \textsuperscript{(20)} میں ایک کیک پک رہا ہے ؛
اُسے اُس نے بچوں کے ہلکے کھانے کے لیے بنایا۔ چولھے کے اوپر
\VUgras{کفگیر} \textsuperscript{(21)} اور \VUgras{جھاگ نکالنے والا
چمچہ} \textsuperscript{(22)} لٹکے ہیں۔

%%K t06-tit-8 sub
\VUsaut{4.0mm}
\VUpk{10.2pt}{40}{برتن۔}
\VUsaut{3.0mm}

%%K t06-c4-06-1 p
6. --- کمرے کے پیچھے لٹکے \VUgras{باورچی خانے کے برتن}
\textsuperscript{(25)} بہت صاف ہیں۔ خوبصورت تانبے کی \VUgras{دیگچیاں}
\textsuperscript{(26)}، \VUgras{دودھ کی ہانڈی} \textsuperscript{(27)}،
\VUgras{چھلنی} \textsuperscript{(28)} اور \VUgras{کدو کش}
\textsuperscript{(29)} احتیاط سے صاف کیے گئے تھے۔

%%K t06-c4-07-1 p
7. --- \VUgras{نمک دان} \textsuperscript{(30)} رکابیوں کی الماری پر
رکھا ہے ؛ اُس کے پاس \VUgras{چائے کا جگ} \textsuperscript{(31)} اور
\VUgras{تیل کی بڑی بوتل} \textsuperscript{(32)} ہیں۔ \VUgras{دراز}
\textsuperscript{(33)} میں شاید باورچی خانے کے چمچے اور چھریاں ہیں۔

%%K t06-c4-08-1 p
8. --- \VUgras{جھاڑو} \textsuperscript{(34)} اور \VUgras{پروں کا
جھاڑن} \textsuperscript{(35)} ایک کونے میں ہیں۔ باورچن نے ابھی
\VUgras{لکڑی کا کندہ} \textsuperscript{(36)} استعمال کیا ؛ اُسے
کھڑکی کے پاس چھوڑ دیا۔

%%K t06-c4-09-1 p
9. --- \VUgras{ہاتھ پونچھنے کا کپڑا} \textsuperscript{(37)} دیوار پر
لٹکا ہے ؛ وہ \VUgras{قیف} \textsuperscript{(38)} کے پاس ہے۔ باورچی
خانے کے اُس حصے میں \VUgras{بھوننے کی جالی} \textsuperscript{(39)}،
\VUgras{قیمہ کرنے کی چھری} \textsuperscript{(40)} اور \VUgras{کاٹنے
کا تختہ} \textsuperscript{(41)} رکھے ہیں۔ پھر، چھری کے اوپر
\VUgras{تختے} \textsuperscript{(42)} پر \VUgras{ہانڈیاں}
\textsuperscript{(43)}، \VUgras{ابالنے کی ہانڈی} \textsuperscript{(44)}
— اُس کا اپنا \VUgras{ڈھکن} \textsuperscript{(45)} ہے — اور
\VUgras{بڑی دیگ} \textsuperscript{(46)} ہیں۔ \VUgras{کڑاہی}
\textsuperscript{(47)} پاس ہی لٹکی ہے۔

%%K t06-c4-10-1 p
10. --- اِس کونے میں صرف تانبے کی \VUgras{ٹونٹی}
\textsuperscript{(48)} چمکتی ہے۔ \VUgras{دھونے کا حوض}
\textsuperscript{(49)} — اُس پر موٹا \VUgras{بڑا گھڑا}
\textsuperscript{(50)} رکھا ہے — اپنی چھوٹی الماری کے ساتھ شاید بہت
آرام دہ ہے ؛ اُس الماری میں \VUgras{سرکے کا پیپا}
\textsuperscript{(51)} — اُس کی اپنی \VUgras{ٹونٹی}
\textsuperscript{(52)} ہے —، \VUgras{کوڑے کا ڈبہ}
\textsuperscript{(53)} اور \VUgras{میلی رکابیاں} \textsuperscript{(54)}
رکھی جاتی ہیں۔

%%K t06-tit-9 sub
\VUsaut{4.0mm}
\VUpk{10.2pt}{40}{باورچی خانے میں رکھی دوسری چیزیں۔}
\VUsaut{3.0mm}

%%K t06-c4-11-1 p
11. --- \VUgras{بڑا ٹب} \textsuperscript{(55)} — اُس میں
\VUgras{سبزیاں} \textsuperscript{(58)} دھوئی جاتی ہیں — اور ڈھالے
لوہے کا \VUgras{کڑھاؤ} \textsuperscript{(56)} — وہ \VUgras{اینٹوں کے
فرش} \textsuperscript{(57)} پر رکھا ہے — شاید اُسی الماری میں اپنی
جگہ رکھتے ہیں۔

%%K t06-c4-12-1 p
12. --- باورچن کو چولھوں کی کمی نہیں ؛ کیونکہ یہ رہا میز کے کونے پر
\VUgras{گیس کا چولھا} \textsuperscript{(59)} ؛ اُس پر چھوٹی دیگچی میں
پانی گرم ہو رہا ہے۔ اُس سے نکلتی \VUgras{بھاپ} \textsuperscript{(60)}
بتاتی ہے کہ وہ ابل رہا ہے۔ یہ پانی کافی کے لیے استعمال ہو گا ؛
کیونکہ یہ رہے میز کے دوسرے سرے پر \VUgras{کافی کا جگ}
\textsuperscript{(64)} اور \VUgras{کافی پیسنے کی چکی}
\textsuperscript{(63)}۔

%%K t06-c4-13-1 p
13. --- چولھا \VUgras{گیس کے جلانے والے} \textsuperscript{(62)} سے
جڑا ہے ؛ لمبی \VUgras{ربڑ کی نلی} \textsuperscript{(61)} اُنھیں
ملاتی ہے۔

%%K t06-c4-14-1 p
14. --- باورچن کی مدد کو آنے والی \VUgras{دہاڑی کی خادمہ}
\textsuperscript{(66)} رات کے کھانے کے لیے سبزیاں تیار کر رہی ہے۔ جب
وہ صاف اور دھل جائیں گی، تو پکانے کا وقت آنے تک اُنھیں \VUgras{بڑے
پیالے} \textsuperscript{(65)} میں رکھے گی۔

%%K t06-tit-10 sub
\VUsaut{4.0mm}
{\centering\fontsize{10.2pt}{10.2pt}\selectfont\textls[40]{\scshape غسل خانہ۔} \textit{(نہانے کا کمرہ۔)}\par}
\VUsaut{3.0mm}

%%K t06-c5-01-1 p
1. --- مکان کے اہم کمرے جاننے کے لیے ہمیں صرف ایک بے تکلفی باقی ہے —
غسل خانے کی ترتیب دیکھنا۔ اُس میں زیادہ وقت نہیں لگے گا، کیونکہ وہ
بڑا نہیں ؛ ہم جلد جان لیں گے کہ \VUgras{غسل کے ٹب}
\textsuperscript{(1)} کے ساتھ اچھا \VUgras{پانی گرم کرنے والا}
\textsuperscript{(2)} ہے، کہ \VUgras{صابن دان} \textsuperscript{(3)}
نہانے والے کے ہاتھ تک پہنچتا ہے، اور کہ \VUgras{تولیہ ٹانگنے کی
ڈنڈی} \textsuperscript{(5)} \VUgras{تولیوں} \textsuperscript{(4)} کو
آسانی سے سکھا دے گی۔

%%K t06-c5-02-1 p
2. --- اِس کمرے کی دیواریں \VUgras{چینی مٹی کی ٹائلوں}
\textsuperscript{(6)} سے ڈھکی ہیں ؛ اُن کی وجہ سے \VUgras{پھوار کا
آلہ} \textsuperscript{(7)} لگانا ممکن ہوا — استعمال کے وقت اچھلنے
والا پانی دیواروں کو خراب کرے گا، اِس کا ڈر نہیں۔ پاؤں دھونے کے کام
آنے والا جست کا \VUgras{چھوٹا حوض} \textsuperscript{(8)} سامان کو
پورا کرتا ہے۔

\VUsaut{6.0mm}
\VUfilet{22mm}
